\label{sec:alg_overview}

Six optimisation algorithms were evaluated across all experimental
configurations, spanning four structural categories: swarm-based,
physics-inspired, ecosystem-inspired, and classical local search.

Particle Swarm Optimization, originally proposed by Kennedy and
Eberhart~\cite{kennedy1995} and extended with an inertia weight by Shi and
Eberhart~\cite{shi1998}, is a swarm method in which each particle maintains
both a personal best position and a shared global best. The inertia weight
$w$ governs the balance between exploration and exploitation through the
velocity update rule. Atomic Orbital Search (AOS)~\cite{azizi2021} is a
physics-inspired method detailed in Section~\ref{sec:algo_select}; it employs
shell-based clustering around a nucleus (global best) as its primary
exploration mechanism, without per-agent personal memory. African Buffalo
Optimization (ABO)~\cite{odili2015} encodes two herd vocalisations --- an
exploratory foraging call and an exploitative regrouping call --- and like
PSO retains both personal and global best positions. Manta Ray Foraging
Optimization (MRFO)~\cite{zhao2020mrfo} models three foraging behaviours
(chain, cyclone, and somersault foraging) but maintains only a global best,
with no per-agent memory. Symbiotic Organisms Search (SOS)~\cite{cheng2014}
simulates mutualistic, commensal, and parasitic ecological interactions,
applying all three phases to every organism per iteration; this yields three
fitness evaluations per organism per iteration and requires no
algorithm-specific hyperparameters. Generalized Pattern Search
(GPS)~\cite{torczon1997} is a classical single-state local optimiser that
evaluates a fixed pattern of directions from the current best solution, with
no population and no historical memory beyond the current position; it serves
as a deliberate baseline to quantify the value of population-based diversity.

Table~\ref{tab:alg_summary} summarises the distinguishing characteristics of
each algorithm.

%TC:ignore
\begin{table}[H]
  \caption{\label{tab:alg_summary}%
    Summary of the six optimisation algorithms evaluated. Memory column
    describes which historical positions are retained per agent. Eval/iter
  denotes fitness evaluations per agent per iteration.}
  \begin{ruledtabular}
    \begin{tabular}{lcllcc}
      Algorithm & Year & Category     & Population-based & Memory
      & Eval/iter \\
      \hline
      AOS       & 2021 & Physics      & Yes              & Global (nucleus)     & 1 \\
      PSO       & 1998 & Swarm        & Yes              & Personal + Global    & 1 \\
      MRFO      & 2020 & Swarm        & Yes              & Global only          & 1 \\
      SOS       & 2014 & Ecosystem    & Yes              & Global only          & 3 \\
      ABO       & 2015 & Swarm        & Yes              & Personal + Global    & 1 \\
      GPS       & ---  & Local search & No               & Current best only    & 1 \\
    \end{tabular}
  \end{ruledtabular}
\end{table}
%TC:endignore
